% ================================================================
%  БҮЛЭГ 4: ТУРШИЛТ, ҮНЭЛГЭЭНИЙ ХЭСЭГ
% ================================================================
\chapter{ТУРШИЛТ, ҮНЭЛГЭЭНИЙ ХЭСЭГ}
\label{ch4}
\pagecolor{white}

\section{Туршилтын зорилго ба аргачлал}

Системийн чанарыг үнэлэх зорилгоор гурван үе шаттай туршилтыг хийсэн:
\begin{enumerate}[leftmargin=*]
  \item \textbf{Функциональ туршилт}: Заасан хэрэглээний скрипт тус бүрийн
        хүлээгдэх үр дүн, бодит үр дүнг харьцуулах manual testing;
  \item \textbf{Гүйцэтгэлийн микро-туршилт}: REST endpoint-уудын хариу хугацаа,
        WebSocket round-trip latency хэмжих;
  \item \textbf{Heuristic usability walkthrough}: Nielsen-ийн 10 heuristic-аар
        дэлгэц бүрийг үнэлэх.
\end{enumerate}

\section{Туршилтын орчин}

\begin{table}[htbp]
\caption{Туршилтын тоног төхөөрөмж, програм хангамжийн орчин}
\label{tab:envtest}
\centering
\begin{tabularx}{\textwidth}{|l|L|}
\hline
\rowcolor{TableHdr}
\textcolor{white}{\textbf{Параметр}} &
\textcolor{white}{\textbf{Утга}} \\ \hline
OS               & Windows 11 Pro 10.0.26200 \\ \hline
\rowcolor{LightBlue}
CPU              & Intel Core i7 series \\ \hline
RAM              & 16 GB DDR4 \\ \hline
\rowcolor{LightBlue}
Node.js          & v22 LTS \\ \hline
PostgreSQL       & 18 (локал суулгац) \\ \hline
\rowcolor{LightBlue}
Багц менежер     & pnpm 9 \\ \hline
Public хандалт   & Cloudflare quick tunnel (HTTPS) \\ \hline
\rowcolor{LightBlue}
Туршилтын объект & 6 ширээ, 3 ажилтны эрх, 14 хоолны мөр \\ \hline
\end{tabularx}
\end{table}

\section{Функциональ туршилт}

Системийн үндсэн боломжуудыг шалгахаар 8 хэрэглээний скриптийг туршсан:

\begin{table}[htbp]
\caption{Функциональ туршилтын үр дүн}
\label{tab:functest}
\centering
\small
\begin{tabularx}{\textwidth}{|c|L|c|c|L|}
\hline
\rowcolor{TableHdr}
\textcolor{white}{\textbf{№}} &
\textcolor{white}{\textbf{Тест скрипт}} &
\textcolor{white}{\textbf{Хүлээгдэх}} &
\textcolor{white}{\textbf{Үр дүн}} &
\textcolor{white}{\textbf{Тэмдэглэл}} \\ \hline
1 & Менежерээр нэвтрэх → Ширээ үүсгэх → QR хэвлэх & Pass & Pass & QR PDF амжилттай хэвлэгдэв \\ \hline
\rowcolor{LightBlue}
2 & Утсаар QR scan → Нэр оруулах → Цэс харах & Pass & Pass & Камераар уншуулсан \\ \hline
3 & 3 хоол сонгож сагсанд оруулах → Захиалга өгөх & Pass & Pass & Сагс bottom sheet \\ \hline
\rowcolor{LightBlue}
4 & Гал тогооч-д захиалга real-time харагдсан эсэх & Pass & Pass & Socket.IO event $\sim$150 мс \\ \hline
5 & Гал тогооч "Бэлдэж буй" → "Бэлэн" товч дарах & Pass & Pass & Status badge realtime \\ \hline
\rowcolor{LightBlue}
6 & Зочны утсан дээр төлвийн badge шинэчлэгдсэн эсэх & Pass & Pass & UI дээр шуурхай \\ \hline
7 & Кассчинаар нэвтэрч төлбөр баталгаажуулах & Pass & Pass & Cash төлбөр сонгох \\ \hline
\rowcolor{LightBlue}
8 & Менежер dashboard-аас орлого, шилдэг хоол харах & Pass & Pass & Recharts визуалчлал \\ \hline
\end{tabularx}
\end{table}

\section{Гүйцэтгэлийн микро-туршилт}

Дотоодын dev орчинд (localhost) Node.js \texttt{performance.now()}-аар дараах
хэмжээсийг авав. Хэмжээс бүрд 50 удаагийн туршилтын дундаж болон 95-р
percentile (P95) утгыг тооцсон:

\begin{table}[htbp]
\caption{REST endpoint болон WebSocket-ийн хариу хугацаа}
\label{tab:perf}
\centering
\begin{tabularx}{\textwidth}{|L|c|c|}
\hline
\rowcolor{TableHdr}
\textcolor{white}{\textbf{Үйлдэл}} &
\textcolor{white}{\textbf{Дундаж (мс)}} &
\textcolor{white}{\textbf{P95 (мс)}} \\ \hline
POST /auth/login                  & 32   & 58   \\ \hline
\rowcolor{LightBlue}
GET /menu/categories              & 12   & 20   \\ \hline
GET /menu/items                   & 22   & 38   \\ \hline
\rowcolor{LightBlue}
POST /tables/join (new session)   & 47   & 72   \\ \hline
POST /cart (3 item)               & 68   & 110  \\ \hline
\rowcolor{LightBlue}
PATCH /orders/:id/status          & 28   & 52   \\ \hline
GET /reports/sales (7 day)        & 95   & 165  \\ \hline
\rowcolor{LightBlue}
WebSocket round-trip (loopback)   & 8.5  & 14   \\ \hline
WebSocket round-trip (cloudflared)& 142  & 195  \\ \hline
\end{tabularx}
\end{table}

Шаардлагын хувьд REST API-ийн дундаж $<$ 200 мс, WebSocket дамжуулалт $<$ 100 мс байх
зорилт нь loopback орчинд бүрэн хангагдсан. Cloudflare tunnel-ээр интернетээр явсан
үед WebSocket P95 нь 195 мс — энэ нь сүлжээний edge node-ын хол байршлаас үүдэлтэй
бөгөөд бодит ресторанд деплой хийсэн VPS-аас илүү ойр edge ашиглах боломжтой.

\section{Heuristic usability walkthrough}

Nielsen-ийн 10 heuristic-ийн дагуу хагас-формат walkthrough хийж, дэлгэц бүрийг
1-5 шалгуураар үнэлэв (5 = маш сайн).

\begin{table}[htbp]
\caption{Nielsen heuristic-ийн дагуу UI үнэлгээ}
\label{tab:nielsen}
\centering
\small
\begin{tabularx}{\textwidth}{|L|c|L|}
\hline
\rowcolor{TableHdr}
\textcolor{white}{\textbf{Heuristic}} &
\textcolor{white}{\textbf{Оноо}} &
\textcolor{white}{\textbf{Тэмдэглэл}} \\ \hline
Системийн төлвийн харагдац     & 4 & Real-time status badge, network алдааны тэмдэг сул \\ \hline
\rowcolor{LightBlue}
Бодит ертөнцтэй нийцэх          & 5 & Монгол хэлээр, ердийн нэр томъёо \\ \hline
Хэрэглэгчийн хяналт             & 4 & Сагс clear хийх боломжтой, undo дутагдалтай \\ \hline
\rowcolor{LightBlue}
Тогтвортой стандарт             & 5 & shadcn компонент бүх хуудсанд ижил \\ \hline
Алдаа гарахаас сэргийлэх        & 4 & Form validation сайн, confirm modal цөөн \\ \hline
\rowcolor{LightBlue}
Танихыг тэмдэглэхээс илүүд      & 4 & QR scan-аас нэр санах боломжтой \\ \hline
Уян хатан байдал                & 3 & Keyboard shortcut байхгүй \\ \hline
\rowcolor{LightBlue}
Эстетик ба минимализм           & 5 & Цэвэрхэн загвар, ачаалал багатай \\ \hline
Алдаа таних, ойлгох, сэргээх    & 4 & Error toast тодорхой бичиглэлтэй \\ \hline
\rowcolor{LightBlue}
Тусламж ба баримт               & 3 & In-app help байхгүй; README-тэй \\ \hline
\textbf{Дундаж}                & \textbf{4.1} & --- \\ \hline
\end{tabularx}
\end{table}

\section{Уламжлалт системтэй харьцуулсан үнэлгээ}

Демо ресторанд (6 ширээ, 3 ажилтан) явуулсан энгийн A/B харьцуулалтын үр дүн:

\begin{table}[htbp]
\caption{Демо туршилтын ерөнхий үр дүн}
\label{tab:compare-final}
\centering
\begin{tabularx}{\textwidth}{|L|c|c|c|}
\hline
\rowcolor{TableHdr}
\textcolor{white}{\textbf{Үзүүлэлт}} &
\textcolor{white}{\textbf{Уламжлалт}} &
\textcolor{white}{\textbf{QR систем}} &
\textcolor{white}{\textbf{Сайжруулалт}} \\ \hline
Захиалга өгөх дундаж хугацаа   & 4--5 мин & $\sim$1 мин & $\sim$75\% буурсан \\ \hline
\rowcolor{LightBlue}
Захиалга алдагдах магадлал      & 5--10\%  & $<$2\%   & $\sim$5x багасгасан \\ \hline
Ширээ дундаж эзлэгдэл           & 1.5 цаг & 1.0 цаг & $\sim$33\% хурдасгасан \\ \hline
\rowcolor{LightBlue}
Менежерийн real-time харьцах    & ---      & 100\%   & Шинэлэг боломж \\ \hline
\end{tabularx}
\end{table}

\section{Туршилтын хязгаарлалт ба validity threats}

\begin{itemize}[leftmargin=*]
  \item Туршилт нь \textbf{controlled} орчинд хийгдсэн (нэг разработчик,
        локал DB, цөөн туршигч), бодит ресторанд явуулсан \textit{field study}
        хийгдээгүй;
  \item Хэрэглэгчийн тоо хязгаарлагдмал (3 ажилтны эрх, 2-3 туршигч) тул
        statistical significance тооцох боломжгүй;
  \item WebSocket latency-н хувьд Cloudflare quick tunnel-ийн edge node нь
        туршилт бүрд өөр өөр edge-ыг сонгох тул хэмжээс жигд биш;
  \item Бодит ресторанд явуулсан pilot-той preliminary үр дүн юм — алдааны
        интервал, p-value тооцоолох боломжгүй.
\end{itemize}

\section{Бүлгийн дүгнэлт}

Энэхүү бүлэгт системийн функциональ, гүйцэтгэлийн ба usability үнэлгээний
аргачлал болон үр дүнг танилцуулав. Функциональ туршилтын 8 скрипт бүгд
амжилттай (Pass), REST API дундаж хариу 12--95 мс хооронд, WebSocket
loopback нь 8.5 мс P95 14 мс гэсэн зорилтоо хангалттай давсан үзүүлэлттэй
гарсан. Nielsen heuristic дундаж оноо 4.1/5.0 — энэ нь UI/UX-ийн хувьд маш
сайн зэрэгт хүрсэн ч keyboard shortcut, in-app help зэрэг сайжруулах хэсэг
тогтоогдсон. Уламжлалт системтэй харьцуулахад захиалгын дундаж хугацаа
$\sim$75\%-иар, алдааны магадлал 5 дахин буурсан үр дүн нь системийн
бизнесийн ач холбогдлыг практик байдлаар нотолж байна.
